%------------------------------------------------------------------------------%
\section{Operações com Séries de Potências}

Apresentamos nessa seção alguns resultados que nos auxiliam a realizar operações
em Séries de Potências. Começamos pela importante propriedade que nos permite
derivar e integrar essas séries termo a termo.
Se uma Série de Potências possui um raio de convergência $R>0$, então a função
\[
  f(x)=\sum_{n=0}^\infty c_n(x-a)^n
\]
está definida no intervalo $(a-R, a+R)$ e é infinitamente derivável, isso é,
tem as derivadas de todas as ordens, neste mesmo intervalo.
Além disso, podemos calcular as derivadas ou integrais de $f$
derivando ou integrando a série termo a termo como descrito nos próximos
dois teoremas

%------------------------------------------------------------------------------%
\begin{teorema}{Derivação de Série de Potências Termo a Termo}{teo:derivacao-termo-termo}
  Se o raio de convergência, $R$, não for nulo, podemos calcular a derivada da função
  \[
    f(x) = \sum_{n=0}^\infty c_n(x-a)^n
  \]
  derivando a série termo a termo dentro do intervalo $(a-R, a+R)$
  \[
    f'(x) \;=\; \sum_{n=o}^\infty \frac{d}{dx}\big( c_n(x-a)^n \big)
          \;=\; \sum_{n=1}^\infty \,n\,c_n\,(x-a)^{n-1}   \;.
  \]
  Essa série também converge em $(a-R, a+R)$.
\end{teorema}
%------------------------------------------------------------------------------%

Observe que o primeiro termo da série é uma constante, portanto quando
derivamos a série ele desaparece e o índice da série para a derivada começa em um.
Note também que podemos aplicar a derivação termo a termo repetidas vezes
para calcular a $k$-ésima derivada de $f$
\begin{align*}
  f^{(k)}(x) & = \sum_{n=0}^\infty \big( c_n(x-a)^n \big)^{(k)}                          \\[3mm]
             & = \sum_{n=k}^\infty \,n\,(n-1) \cdots \big(n-(k-1)\big)\,c_n\,(x-a)^{n-k} \;.
\end{align*}

%------------------------------------------------------------------------------%
\begin{teorema}{Integração de Série de Potências Termo a Termo}{teo:integracao-termo-termo}
  Se o raio de convergência, $R$, não for nulo, podemos calcular a primitiva da função
  \[
    f(x) = \sum_{n=0}^\infty c_n(x-a)^n
  \]
  integrando a série termo a termo dentro do intervalo $(a-R, a+R)$
  \[
    \int f(x)\,dx \;=\; \sum_{n=0}^\infty \int c_n(x-a)^n dx
                  \;=\; \sum_{n=0}^\infty c_n\frac{(x-a)^{n+1}}{n+1}+ C
  \]
  onde $C$ é a constante de integração. Esta série também converge em $(a-R, a+R)$.
\end{teorema}
%------------------------------------------------------------------------------%

Esses teoremas, combinados com a fórmula que temos para a Série Geométrica,
são úteis para encontrar a função que representa uma dada série,
ou para encontrar uma série para uma dada função. Claro que isso só será possível quando
operações de integração ou derivação nos dão ou uma série geométrica ou uma função que
pode ser vista como limite de uma série geométrica. Os próximos exemplos ilustram essa
aplicação.

%------------------------------------------------------------------------------%
\begin{example}
  Encontre uma série de potências centrada em $x=0$ para a função
  \(
    f(x)=\ln \left( 1+x^2 \right)
  \).
  \solution
  Note que
  \[
    \frac{d}{dx} \ln\left( 1+x^2 \right)
      = \frac{2x}{1+x^2}
      = \frac{2x}{1-\left(-x^2\right)}   \;.
  \]
  Comparando essa expressão com a fórmula da soma
  Série Geométrica com $r=-x^2$
  \[
    \sum_{n=0}^\infty \left(-x^2\right)^n=\frac{1}{1-\left(-x^2\right)}
  \]
  percebemos que, nos pontos onde a série converge,
  \[
    \frac{d}{dx}\ln\left( 1+x^2 \right)
      = 2x \sum_{n=0}^\infty \left(-x^2 \right)^n
      = 2  \sum_{n=0}^\infty (-1)^n\,x^{2n+1}     \;.
  \]
  Integrando e simplificando temos
  \[
    \ln\left( 1+x^2 \right) = \sum_{n=0}^\infty (-1)^n\frac{x^{2n+2}}{\,n+1\,}+C
  \]
  onde $C$ é uma constante.
  Avaliando a série em $x=0$ verificamos que $C=\ln(1)=0$, portanto
  \[
    \ln\left( 1+x^2 \right) = \sum_{n=0}^\infty (-1)^n\frac{x^{2n+2}}{n+1}     \;.
  \]
  Note que provamos a validade dessa fórmula apenas para $x\in(-1,1)$,
  pois usamos o limite da Série Geométrica que converge apenas para valores
  de $x$ tais que
  \(
    \abs{r}=\abs{-x^2}<1     \;.
  \).
  Porém, série que aparece na última equação
  converge no intervalo $[-1, 1]$, entretanto, neste momento
  só podemos afirmar que a igualdade vale para $x\in (-1, 1)$.
  Os casos $x=-1$ e $x=1$ precisariam ser analisados separadamente.
\end{example}
%------------------------------------------------------------------------------%

No próximo exemplo faremos o caminho inverso do anterior, ou seja,
começaremos com uma série e procuraremos uma função que represente
essa série em algum intervalo.

%------------------------------------------------------------------------------%
\begin{example}
  Encontre uma função $f(x)$ que represente a série
  \[
    \sum_{n=0}^\infty (-1)^n\frac{x^{2n+1}}{\,2n+1\,}
  \]
  em um intervalo contendo a origem.

  \solution
  Se
  \[
    f(x) = \sum_{n=0}^\infty (-1)^n\frac{x^{2n+1}}{\,2n+1\,}
  \]
  podemos derivar $f$ termo a termo
  \[
    f'(x) = \sum_{n=0}^\infty (-1)^n \frac{d}{dx}\left(\frac{x^{2n+1}}{\,2n+1\,}\right)
          = \sum_{n=0}^\infty (-1)^nx^{2n}
          = \sum_{n=0}^\infty \left(-x^2 \right)^n   \;.
  \]
  Rearranjando temos
  \[
    f'(x) = \sum_{n=1}^\infty \left(-x^2 \right)^{n-1}   \;.
  \]
  Comparando com uma Série Geométrica observamos que $a=1$ e $r=-x^2$.
  Usando a fórmula para a soma da série temos que,
  para $x\in (-1, 1)$,
  \[
    f'(x) = \frac{1}{\,1-\left(-x^2\right)\,}
          = \frac{1}{\,1+x^2\,}     \;.
  \]
  Veja que $f(0)=0, $ então integrando $f'(t)$ sobre $[0, x]$
  e usando o Teorema Fundamental do Cálculo~\ref{teo:fundamental-calculo},
  obtemos que
  \begin{align*}
    f(x) & = \int_0^x f'(t)\,dt               \\[3mm]
         & = \int_0^x \frac{1}{\,1+t^2\,}\,dt \\[3mm]
         & = \arctan(x) - \arctan(0)          \\[3mm]
         & = \arctan(x)     \;.
  \end{align*}
  Note que só podemos garantir que a função $\arctan(x)$ representa a
  série no intervalo $(-1, 1)$, pois utilizamos o limite da Série Geométrica,
  que só converge neste intervalo.
  Observe ainda que podemos afirmar que essa série converge no extremo $x=1$,
  usando o Teste de Leibniz~\ref{teo:Teste-Leibniz}.
  Entretanto, neste momento, não podemos afirmar que
  em $x=1$ essa série converge para $\arctan(1)$.
\end{example}
%------------------------------------------------------------------------------%

Apresentamos a seguir algumas operações que podem ser realizadas em séries de
potências dentro do seu intervalo de convergência.

%------------------------------------------------------------------------------%
\begin{teorema}{Multiplicação de Séries de Potências}{teo:multipliacao-series-potencias}
	Se as séries
	\[
		A(x) = \sum_{n=0}^{\infty} a_n x^n
    \qquad\text{ e }\qquad
  	B(x) = \sum_{n=0}^{\infty} b_n x^n
	\]
	convergem absolutamente para \(\abs{x}<R\), e
	\[
		c_n = \sum_{k=0}^{n} a_kb_{n-k}
         = a_0b_n + a_1b_{n-1} + a_2b_{n-2} + \cdots + a_{n-1}b_1 + a_nb_0     \;,
	\]
	então
  \[
    \sum_{n=0}^{\infty} c_nx^n = A(x)B(x)
    \qquad\qquad \abs{x}<R     \;,
  \]
	ou seja, podemos escrever
	\[
		\left(\sum_{n=0}^{\infty} a_n x^n\right)
		\left(\sum_{n=0}^{\infty} b_n x^n\right) =
		\sum_{n=0}^{\infty} c_nx^n     \;.
	\]
\end{teorema}
%------------------------------------------------------------------------------%

Observe que esse resultado nos diz que o produto de duas séries de potencia convergentes
também converge no mesmo intervalo. Porém, calcular os valores $c_n$ pode ser uma
tarefa extremamente complexa.

Podemos também fazer uma mudança de variáveis trocando $x$ por uma função $f(x)$.

%------------------------------------------------------------------------------%
\begin{teorema}{Mudança de Variáveis em uma Série de Potências}{teo:Mudanca-variaveis-serie-potencias}
	Se a série
	\[
	  \sum_{n=0}^{\infty} a_n x^n
	\]
	converge absolutamente para \(\abs{x}<R\), então
	\[
  	\sum_{n=0}^{\infty} a_n \big(\, f(x) \,\big)^n
	\]
	converge absolutamente desde que $f$ seja uma função contínua e
	\(
    \abs{f(x)}<R
  \).
\end{teorema}
%------------------------------------------------------------------------------%

%%------------------------------------------------------------------------------%
%\begin{teorema}{Composição de Funções}{teo:composicao-funcoes}
%  Seja $f$ uma função satisfazendo
%  \[
%    f(x) = \sum_{n=0}^\infty c_n(x-a)^n \qquad \text{ em } x\in I
%  \]
%  onde $I$ é um intervalo aberto.
%  Seja $u:\R\rightarrow \R$ uma função contínua e $(b, c)$ é o intervalo $u^{-1}(I)$. Então,
%  \[
%    f\big(u(x)\big) = \sum_{n=0}^\infty c_n \big(u(x)-a\big)^n \qquad \text{ em } x\in (b, c)
%  \]
%\end{teorema}
%%------------------------------------------------------------------------------%
%\begin{proof}
%  Como $f$ é uma série de potências, sabemos, pela Proposição \ref{prop:serietaylorpotencia}
%  que $f$ coincide com sua série de Taylor em $I$. Então, pelo teorema de Taylor temos que
%  \[
%    f(x) = \sum_{n=0}^k \frac{f^{(n)}(a)}{n!} (x-a)^n+R_{k}(x)
%  \]
%  e que $R_k(x)\rightarrow 0$ para todo $x\in I$. Então,
%  \[
%    f\big(u(x)\big) = \sum_{n=0}^k \frac{f^{(n)}(a)}{n!} \big(u(x)-a\big)^n + R_{k}(u(x))
%  \]
%  e $R_{k}\big(u(x)\big)\rightarrow 0$ se $u(x)\in I$,
%  mas se $x\in (b, c)=u^{-1}(I)$ segue que $u(x)\in I$.
%  Portanto, sempre que $x\in (b, c)$,
%  \[
%    f\big(u(x)\big) = \sum_{n=0}^\infty \frac{f^{(n)}(a)}{n!} \big(u(x)-a\big)^n
%    = \sum_{n=0}^\infty c_n \big(u(x)-a\big)^n
%  \]
%\end{proof}
%%------------------------------------------------------------------------------%

%------------------------------------------------------------------------------%
